% 版式、页眉页脚、定理编号 
% 行距、段距、标题字号、页面风格
% ===== 定理环境 =====
\theoremstyle{definition}
\newtheorem{theorem}{定理}[section]
\newtheorem{proposition}{命题}[section]
\newtheorem{lemma}{引理}[section]
\newtheorem{corollary}{推论}[section]
\newtheorem{definition}{定义}[section]
\newtheorem{example}{例}[section]

% ===== 段落样式 =====
% \setlength{\parindent}{2em}
% \setlength{\parskip}{0.5em}

% ===== 行距 =====
\linespread{1.3}

% ===== 段落间距 =====
\setlength{\parskip}{0.6em}

% ===== 段首不缩进（讲义必选）=====
\setlength{\parindent}{0pt}
% =========================
% 颜色定义（仅 ASCII）
% =========================
\usepackage{xcolor}

% \definecolor{CStatement}{RGB}{52, 152, 219}   % 蓝
% \definecolor{CExplain}{RGB}{46, 204, 113}     % 绿
% \definecolor{CProof}{RGB}{155, 89, 182}       % 紫
% \definecolor{CExample}{RGB}{241, 196, 15}     % 黄
% \definecolor{CTrap}{RGB}{231, 76, 60}         % 红
% \definecolor{CSummary}{RGB}{127, 140, 141}    % 灰

% ========= 数学安全宏 =========
\newcommand{\dd}{\,\mathrm{d}}
\newcommand{\leqsl}{\leqslant}
\newcommand{\geqsl}{\geqslant}

\usepackage{etoolbox}
\pretocmd{\subsubsection}{\clearpage\phantomsection}{}{}

% =========================
% Conclusion page break helper
% - No sectioning inside
% - Only: new page + correct hyperlink anchor
% =========================
\newcommand{\Conclusion}{%
  \clearpage
  \phantomsection
}

% =========================
% Header / Footer (公众号)
% =========================
\usepackage{fancyhdr}

% 公众号样式颜色（推荐：暗红）
\definecolor{wechatred}{RGB}{140,20,20}
\definecolor{brandblue}{RGB}{0, 75, 150}
\definecolor{darkgray}{RGB}{80, 80, 80}

\pagestyle{fancy}
\fancyhf{} % 清空默认

\setlength{\headheight}{14pt} %增大：给页眉留足空间，不再报警
\addtolength{\topmargin}{-2pt} %往上挪一点：避免正文被挤下去导致版心变差


% 页眉：左边显示章节名，右边显示小节名
% \fancyhead[L]{\small\bfseries\leftmark}
% \fancyhead[C]{\small 加V: \WX 免费领取配套例题与更新 |本文档免费}
% \fancyhead[R]{\small\rightmark}

\fancyhead[L]{\small\bfseries\color{gray} 高中数学 · 提分利器 \big| \color{black} 必修系列}
\fancyhead[R]{\small \rightmark}
\renewcommand{\headrulewidth}{0.5pt} % 增加一条细线提升质感

% 页脚：左边公众号提示，右边页码
% \fancyfoot[L]{\small\color{wechatred}{\textbf{关注公众号:}\,\WX 打造最全数学二级结论库 |加V: \WX 免费领取配套例题与更新 |本文档免费}}
% \fancyfoot[R]{\small\thepage}

% 页脚设置
\fancyfoot[L]{
    \small\color{darkgray} 
    \textbf{高中数学·秒杀系列} | 内部教研资料 \hfill \\
    \scriptsize \textit{注：由 \WX 整理分享，旨在助力学子高效提分}
}

\fancyfoot[C]{} % 中间留空

\fancyfoot[R]{
    \small 
    \color{wechatred}\textbf{更新见公众号：} \textbf{\WX} \\
    \color{black}\footnotesize 第 \thepage \ 页 \ / \ 共 \pageref{LastPage} 页
}

% 调整页脚横线（可选）
\renewcommand{\footrulewidth}{0.4pt}

% 页眉页脚分隔线
% \renewcommand{\headrulewidth}{0.4pt}
\renewcommand{\footrulewidth}{0.4pt}

% 让章节标题正确写入页眉
\renewcommand{\sectionmark}[1]{\markboth{#1}{}}
\renewcommand{\subsectionmark}[1]{\markright{#1}}

%==============================
% Watermark (讲义水印设置)
%==============================
% \usepackage{draftwatermark}
% \SetWatermarkText{远锋数学讲义 · 禁止转载}
% \SetWatermarkScale{3.6}
% \SetWatermarkAngle{45}
% \SetWatermarkLightness{0.93} % 越接近1越淡, 0.90~0.95最好


%================================================
% Watermark (no header/footer modification)
%================================================

\usepackage{tikz}
\usepackage{eso-pic}

% 你只需要改这2行
\newcommand{\WMHiddenID}{ID: 2026E03-001}


%--- 隐藏ID（右上角极淡）
\newcommand{\LectureHiddenID}{%
  \AddToShipoutPictureFG*{%
    \begin{tikzpicture}[remember picture,overlay]
      \node[
        anchor=north east,
        text=black,
        opacity=0.12,
        scale=0.85
      ] at ([xshift=-8mm,yshift=-6mm]current page.north east)
      {\WMHiddenID};
    \end{tikzpicture}%
  }%
}

%--- 自动启用（不影响你的页眉页脚）
\AtBeginDocument{%
  \LectureHiddenID
}


% =========================
% 🚀 不可裁剪底边标识（关键增强）
% =========================
\usepackage{eso-pic}

\AddToShipoutPicture{
\begin{tikzpicture}[remember picture, overlay]

% 底部贴边（几乎贴到物理边界）
\node[
    anchor=south,
    yshift=2pt, % 非常贴边
    opacity=0.9
] at (current page.south)
{
    \scriptsize
    © Math \& Code Lab ｜公众号：\WX ｜内部教研资料
};

\end{tikzpicture}
}

% =========================
% 🚀 左右边界锁定（防横向裁剪）
% =========================
\AddToShipoutPicture{
\begin{tikzpicture}[remember picture, overlay]

% 左边
\node[
    rotate=90,
    anchor=south,
    xshift=3pt,
    opacity=0.15
] at (current page.west)
{\tiny Math \& Code Lab};

% 右边
\node[
    rotate=-90,
    anchor=south,
    xshift=-3pt,
    opacity=0.15
] at (current page.east)
{\tiny \WX};

\end{tikzpicture}
}

% =========================
% 🚀 顶部补刀（防极端裁剪）
% =========================
\AddToShipoutPicture{
\begin{tikzpicture}[remember picture, overlay]

\node[
    anchor=north,
    yshift=-2pt,
    opacity=0.2
] at (current page.north)
{
    \scriptsize 关注公众号 \WX, 免费领取更多资料
};

\end{tikzpicture}
}

% =========================
% 🚀 真·防搜索版本（宏拆分 + 断字符）
% =========================

\newcommand{\wxA}{ok}
\newcommand{\wxB}{-}
\newcommand{\wxC}{shuxue}

% =========================
% 🚀 真·防搜索（Unicode级）
% =========================

\newcommand{\zwsp}{\char"200B} % zero width space

\newcommand{\WX}{
o\zwsp k\zwsp-\zwsp s\zwsp h\zwsp u\zwsp x\zwsp u\zwsp e
}

% 定义带圈数字命令 \circled{1}
\newcommand{\circled}[1]{%
  \begin{tikzpicture}[baseline=(char.base)]
    \node[shape=circle,draw,inner sep=1pt, minimum size=1.2em] (char) {\small #1};
  \end{tikzpicture}%
}